\chapter{1986:Gerd Faltings}\label{ch1986gerdfaltings:gerd-faltings-fields-medal-1986-a-comprehensive-report}

\begin{center}\rule{0.5\linewidth}{0.5pt}\end{center}

\textbf{Prepared:} July 2026 \textbf{Subject:} Gerd Faltings (b.~July 28, 1954, Gelsenkirchen-Buer, West Germany) \textbf{Position:} Professor, University of Bonn; Director, Max Planck Institute for Mathematics (since 1995) \textbf{Doctoral Advisor:} Hans-Joachim Nastold (University of M\"unster) \textbf{Fields Medal Citation:} \emph{``For his proof of the Mordell conjecture, thereby establishing a fundamental result in Diophantine geometry.''}

\begin{center}\rule{0.5\linewidth}{0.5pt}\end{center}

\begin{prereqs} Algebraic geometry (schemes, abelian varieties, moduli spaces), number theory (elliptic curves, Galois representations), basic homological algebra and Arakelov theory. \end{prereqs}

\begin{objectives} Explain the Mordell conjecture and its significance in Diophantine geometry; describe Faltings's proof strategy via the Shafarevich and Tate conjectures; understand the Faltings height and its role in finiteness arguments. \end{objectives}

\section{Problem History}

The Mordell conjecture, formulated by Louis Mordell in 1922, asserts that a smooth projective curve of genus $g \geq 2$ over a number field $K$ has only finitely many $K$-rational points. This was one of the central open problems in Diophantine geometry for over six decades. The conjecture expressed a deep intuition: geometric complexity, measured by genus, should control arithmetic complexity, measured by the number of rational points.

For curves of genus $0$ (the projective line), rational points are infinite over any number field, parameterized by the line itself. For curves of genus $1$ (elliptic curves), rational points can be finite or infinite, and the Birch--Swinnerton-Dyer conjecture predicts how to determine the rank. For genus $g \geq 2$, the Mordell conjecture asserts a sharp dichotomy: the arithmetic becomes fundamentally more constrained.

Partial results preceded Faltings's work. Siegel proved in 1929 that integral points on curves of genus $\geq 1$ are finite, but this is a different problem---it concerns points whose coordinates are algebraic integers, not merely algebraic numbers. Chabauty developed a $p$-adic method that proves finiteness when the rank of the Jacobian is less than the genus, but this condition fails in general. Mumford showed using $p$-adic methods that rational points on abelian varieties lie in finitely many cosets of subgroups, but this did not directly resolve the conjecture.

The Mordell conjecture was considered extremely difficult. Many mathematicians attempted proofs over the years, and the problem became a touchstone for the development of new techniques in arithmetic geometry. The failure of elementary methods led to the recognition that fundamentally new ideas---connecting moduli spaces, heights, and Galois representations---would be required.

A related conjecture, formulated by Shafarevich in the early 1960s, predicted that there are only finitely many isomorphism classes of abelian varieties of dimension $g$ over $K$ with good reduction outside a given finite set of primes $S$. This finiteness question, while seemingly more technical, turned out to be the key ingredient in the Mordell conjecture.

The Tate conjecture for abelian varieties over number fields---predicting that the Galois representation on the Tate module completely determines the isogeny class---was another major open problem. Faltings would prove both the Shafarevich and Tate conjectures as intermediate steps, making his proof a landmark that resolved three conjectures simultaneously.

The historical context of the problem is rooted in the arithmetic geometry tradition stretching from Diophantus through Fermat, Euler, and Mordell. The conjecture connects to the broader Langlands program, to the theory of heights, and to the study of rational points on higher-dimensional varieties via the Lang conjectures.

\section{How Was It Solved}

Faltings announced his proof at a conference in Oberwolfach in 1983, and published the 18-page paper ``Endlichkeitss\"atze f\"ur abelsche Variet\"aten \"uber Zahlk\"orpern'' in \emph{Inventiones Mathematicae} in 1984. The proof is a masterpiece of compression, establishing several deep results in rapid succession.

The proof proceeds in four major steps. First, Faltings proved the Shafarevich conjecture: given $K$, $g$, and a finite set $S$, there are finitely many isomorphism classes of $g$-dimensional abelian varieties over $K$ with good reduction outside $S$. Second, he proved the Tate conjecture for abelian varieties over number fields. Third, he established an isogeny bound. Fourth, he combined these to deduce finiteness of rational points.

The proof of the Shafarevich conjecture uses the theory of moduli spaces of abelian varieties. Abelian varieties with good reduction outside $S$ correspond to integral points on the Siegel modular variety $\mathcal{A}_{g}$ (with appropriate level structure). The Faltings height provides a bound on these points: it is an Arakelov-theoretic invariant that measures the ``arithmetic complexity'' of an abelian variety, defined using the N\'eron model and the sheaf of invariant differential forms:
\[
h_F(A) = \frac{1}{[K:\mathbb{Q}]} \widehat{\deg}(\omega_{\mathcal{A}/\mathcal{O}_K})
\]
where $\widehat{\deg}$ denotes the Arakelov degree incorporating contributions from both finite and archimedean places. The key property is that there are only finitely many isomorphism classes of abelian varieties with bounded Faltings height.

The proof of the Tate conjecture uses the semisimplicity of the Tate module as a Galois representation. For two abelian varieties $A$ and $B$ over $K$, the natural map
\[
\operatorname{Hom}_K(A, B) \otimes \mathbb{Q}_l \to \operatorname{Hom}_{G_K}(T_l(A), T_l(B))
\]
is an isomorphism, meaning the Galois representation on the Tate module $T_l(A) = \varprojlim_n A[l^n]$ completely determines the isogeny class. This is essential for connecting the set of rational points on a curve to abelian varieties with prescribed properties.

The final step uses the Jacobian variety $J(C)$ of a curve $C$ of genus $g \geq 2$. Rational points on $C$ map to $J(C)$ via the Abel--Jacobi map. If $C$ had infinitely many rational points, their images would give infinitely many points with properties that contradict the finiteness from the Shafarevich conjecture, using the Tate conjecture to control the isogeny classes.

The proof introduced several technical innovations: Arakelov intersection theory on the N\'eron model, the theory of heights on moduli spaces, the construction of appropriate integral models of Siegel modular varieties, and the use of toroidal compactifications to handle boundary behavior. Each of these tools has since become standard in arithmetic geometry.

A crucial ingredient is the theory of N\'eron models. Given an abelian variety $A$ over the ring of integers $\mathcal{O}_{K,S}$ of $K$ with good reduction outside $S$, the N\'eron model provides a smooth group scheme over $\mathcal{O}_{K,S}$ whose generic fiber is $A$. The height function on the N\'eron model measures the arithmetic complexity of the reduction, and Faltings showed that this height can be bounded using Arakelov-theoretic methods.

The Arakelov degree of a metrized line bundle $\overline{L}$ on an arithmetic variety $X$ over $\mathcal{O}_K$ is defined as:
\[
\widehat{\deg}(\overline{L}) = \log \#(L / s \cdot \mathcal{O}_X) + \sum_{v \mid \infty} \log \|s\|_v
\]
where the first term is the contribution from finite places and the sum is over archimedean places. This degree is fundamental to the definition of the Faltings height and underlies the entire framework of Arakelov geometry.

The isogeny bound established by Faltings states that the number of isogenies between abelian varieties of bounded height and dimension is finite. This follows from a deep analysis of the distribution of isogenies on moduli spaces, using the theory of heights and the compactness of the moduli space of abelian varieties.

The role of the Jacobian in the proof deserves further emphasis. For a curve $C$ of genus $g \geq 2$, the Jacobian $J(C)$ is a $g$-dimensional abelian variety, and the Abel--Jacobi map $\kappa: C \to J(C)$ sends a base point $p_0$ and a point $p$ to the class of $p - p_0$ in the Jacobian. The image $\kappa(C)$ spans $J(C)$ as a group, and the induced map on rational points is injective for curves of genus $\geq 2$ (since the class of a single point cannot be torsion). Thus the finiteness of $C(K)$ follows from the finiteness of rational points on $\kappa(C) \subset J(C)$ with prescribed Galois-theoretic properties.

The construction of integral models of Siegel modular varieties is another key technical ingredient. The Siegel modular variety $\mathcal{A}_g$ parameterizes $g$-dimensional abelian varieties with a chosen level structure. Faltings constructed appropriate integral models over $\mathcal{O}_K$ that are smooth away from a normal crossings divisor, and used Arakelov intersection theory on these models to bound the Faltings height. The construction involves choices of level structures that are compatible with the reduction theory of abelian varieties.

\section{Impact of This Discovery}

Faltings's proof immediately transformed Diophantine geometry. The concept of the Faltings height became a central tool, used subsequently in the study of heights on Shimura varieties, the Bogomolov conjecture, and the distribution of rational points. The methods of the proof---especially the interplay between moduli spaces, Arakelov theory, and Galois representations---influenced nearly every subsequent advance in the field.

The proof of the Tate conjecture had consequences extending well beyond the Mordell conjecture. It established that the structure of the Tate module as a Galois representation completely determines the isogeny class, providing a bridge between the arithmetic and the representation-theoretic aspects of abelian varieties. This is closely related to the Mumford--Tate conjecture, which predicts a similar relationship between the Hodge structure and the Galois representation.

The resolution of the Shafarevich conjecture provided a new perspective on the finiteness of abelian varieties with prescribed reduction properties. This result is independently important, connecting to the theory of Shimura varieties, the construction of modular forms, and the study of arithmetic of abelian varieties over local and global fields.

Beyond the specific conjectures proved, Faltings's work demonstrated the power of Arakelov theory in arithmetic geometry. He developed an arithmetic Riemann--Roch theorem:
\[
\widehat{\operatorname{ch}}(f_* \mathcal{E}) = f_*(\widehat{\operatorname{ch}}(\mathcal{E}) \cdot \widehat{\operatorname{Td}}(T_f))
\]
relating the arithmetic Chern character to the pushforward, which has become a fundamental tool for computing heights and invariants on arithmetic varieties.

The proof also had profound implications for the broader Langlands program. Faltings's work on the Fontaine--Faltings comparison theorem relates $p$-adic \'etale cohomology to de Rham cohomology via the period ring $\mathbb{B}_{\mathrm{dR}}$, establishing that $p$-adic representations arising from algebraic geometry are ``geometric'' in the sense of Fontaine. This is a key ingredient for the Fontaine--Mazur conjecture.

Faltings extended his methods to prove the Hodge conjecture for abelian varieties of CM type and certain Shimura varieties, and he proved a generalization of the Mordell conjecture to higher-dimensional varieties with Bombieri, following Vojta's approach. The Faltings height continues to find new applications, from the Gross--Zagier formula to the study of special points on Shimura varieties.

The Chowla--Selberg formula, generalized by Faltings, expresses the Faltings height of a CM abelian variety in terms of values of the Gamma function at algebraic arguments:
\[
h_F(A) = \sum_{v \mid \infty} \log \left( \prod_{i=1}^{g} \Gamma\left(\frac{k_i}{n}\right) \right) + \text{algebraic terms}
\]
where the $k_i$ are determined by the CM type. This formula has deep connections to modular forms and special values of $L$-functions.

\section{Modern Views and Improvements}

The most dramatic extension of Faltings's work came from the proof of the finite generation of the canonical ring by Birkar, Cascini, Hacon, and McKernan in 2010, which completed the minimal model program and provided the deepest understanding of birational classification. This work, which earned Caucher Birkar a Fields Medal in 2018, builds on the same moduli-theoretic foundations that Faltings used.

The Bogomolov conjecture, proved by Ullmo (1998) and Zhang (1998) using ideas closely related to Faltings's height theory, asserts that algebraic points of small height on a curve of genus $g \geq 2$ are close to a finite set of special points. This is a metric analogue of the Mordell conjecture, replacing the finiteness of rational points with the equidistribution of points of small height.

The theory of $p$-adic Hodge theory, foundations of which were laid by Faltings, has been revolutionized by Peter Scholze's theory of perfectoid spaces (Fields Medal 2018). Scholze's almost purity theorem generalizes Faltings's result in a profound way, using perfectoid algebras to establish new comparison theorems between different cohomology theories.

The Langlands program has continued to evolve, with the geometric Langlands program (formulated by Drinfeld and Beilinson) providing a geometric framework that connects Faltings's work to quantum field theory via S-duality. The relationship between the Faltings height and special values of $L$-functions, anticipated in Faltings's work, remains an active area of research.

Vojta's dictionary, which provides a dictionary between Diophantine approximation and Nevanlinna theory, gives a broader perspective on Faltings's results. The conjectures of Lang and Vojta generalize the Mordell conjecture to higher-dimensional varieties, and Faltings's methods provide tools that continue to be developed in this direction.

The Chowla--Selberg formula, generalized by Faltings to higher-dimensional abelian varieties of CM type, connects the Faltings height to special values of the Gamma function. This formula has been used in the study of special points on Shimura varieties and in the arithmetic of modular curves.

Modern computational methods have also made the Mordell conjecture more concrete. Algorithms for computing rational points on curves of genus $\geq 2$, based on Chabauty--Coleman methods and Faltings's height theory, have been implemented and can now handle curves of significant genus and conductor. The Magma and SageMath computer algebra systems include routines for computing Faltings heights and testing finiteness of rational points.

The proof of the full Langlands functoriality conjecture for functorial lifts, by Arthur and others, extends the framework that Faltings contributed to. The trace formula methods and automorphic representation theory that Faltings's results connect to continue to be central tools in the Langlands program.

The Faltings height has found unexpected applications in the theory of special values of $L$-functions. The Deligne conjecture on special values of $L$-functions, now largely proved, predicts that $L$-values can be expressed in terms of periods and algebraic invariants. The Faltings height appears naturally in these expressions for abelian varieties, connecting the arithmetic of heights to the analytic theory of $L$-functions.

The generalization of Faltings's theorem to higher-dimensional varieties remains an active area of research. The Lang conjecture predicts that for a variety $X$ of general type over a number field $K$, the set of rational points $X(K)$ is not Zariski dense. Faltings and Bombieri proved special cases of this conjecture, but the general case is far from resolved. The methods of Faltings---heights, moduli spaces, and Arakelov theory---continue to be the main tools for attacking this problem.

The Vojta conjecture, which generalizes the Mordell conjecture to higher-dimensional varieties and subsumes many other conjectures in Diophantine geometry, was the subject of intense study following Faltings's proof. Vojta's approach, using Nevanlinna-theoretic methods and the theory of heights, provided an alternative perspective on Faltings's results. The dictionary between Diophantine approximation and Nevanlinna theory, anticipated by Vojta, has been a productive framework for understanding the arithmetic of algebraic varieties.

The relationship between the Faltings height and the stable Faltings height---defined using the moduli space of abelian varieties---has been clarified by many researchers. The stable Faltings height is an invariant of the isogeny class of an abelian variety, and its distribution on moduli spaces is governed by the Siegel metric. The study of these heights has applications to the arithmetic of modular forms, the distribution of $L$-values, and the theory of special points on Shimura varieties.

The Mordell--Weil theorem, which states that the group of rational points on an abelian variety is finitely generated, is a precursor to Faltings's work. The Mordell--Weil theorem guarantees that $J(K)$ is finitely generated, but it does not say that $C(K)$ is finite. Faltings's theorem goes further by showing that the image of $C(K)$ in $J(K)$ under the Abel--Jacobi map is finite, using the fact that the genus condition forces the image to be contained in a proper abelian subvariety of $J(K)$ for large genus.

The Tate conjecture, proved by Faltings, has implications for the study of zeta functions of abelian varieties. The zeta function of an abelian variety $A$ over $\mathbb{F}_q$ encodes the number of points on $A$ over finite extensions of $\mathbb{F}_q$, and the Tate conjecture determines the structure of the Galois representation on the Tate module from the zeta function. This connection between the arithmetic of abelian varieties and the analytic theory of zeta functions is a key theme in modern number theory.

The Faltings height has been used to study the distribution of abelian varieties over number fields. The distribution of abelian varieties by height is governed by the Siegel measure on the moduli space, and the Faltings height provides a natural height function for this distribution. The study of this distribution has applications to the arithmetic of modular forms, the distribution of $L$-values, and the theory of special points on Shimura varieties.

The connection between Faltings's work and the theory of special cycles on Shimura varieties is an active area of research. The special cycles on Shimura varieties are subvarieties that correspond to abelian varieties with additional structure, and the intersection theory of these cycles is related to the Faltings height and the Gross--Zagier formula. The study of special cycles has applications to the Langlands program and the theory of automorphic forms.

\begin{center}\rule{0.5\linewidth}{0.5pt}\end{center}

\section*{Summary}
\addcontentsline{toc}{section}{Summary}

Faltings proved the Mordell conjecture (now Faltings's theorem): every curve of genus $g \geq 2$ over a number field has only finitely many rational points. His proof established the Shafarevich and Tate conjectures for abelian varieties, introduced the Faltings height, and laid foundations for $p$-adic Hodge theory and Arakelov geometry that continue to shape arithmetic geometry.
